% GENERATED FILE. Edit canonical lessons, then run npm run generate:books.
% canonical-source-hash: fnv1a64:8df32eccae4b99c2
% canonical-lessons: RU-C364-obshchestvo, RU-C364-naselenie, RU-C364-obshchestvennyy, RU-C364-socialnyy, RU-C364-zhitel

\chapter{Society and its people}
\label{ch:ru-society-and-its-people}
\hlchaptermodality{364}

\begin{chapteropening}
\textbf{By the end of this chapter:} \emph{I can talk about society, the population of a place and its inhabitants, and about public and social questions.}

Complete the last lesson of chapter 364: I can talk about society, the population of a place and its inhabitants, and about public and social questions.
\end{chapteropening}

\section[\texorpdfstring{\ru{общество} (óbshchestvo)}{общество (óbshchestvo)}]{\ru{общество} (óbshchestvo) — society}

\label{lesson:RU-C364-obshchestvo}



\begin{quote}
Before the new one: say the Russian for a competitor, then the Russian for to cooperate.
\end{quote}

\subsection*{You'll want to know: \ru{общество}}
\textbf{\ru{общество}} (\emph{óbshchestvo}) — \textquotedblleft{}society\textquotedblright{}.

\begin{quote}
\ru{В} \ru{современном} \ru{обществе} \ru{много} \ru{проблем}.
\end{quote}

\begin{quote}
\emph{V sovreménnom óbshchestve mnógo problém.}
\end{quote}

\begin{quote}
There are many problems in modern society.
\end{quote}

Neuter. It also means a society or club: \emph{\ru{общество} \ru{любителей} \ru{книги}}.

\begin{cousinweb}
\textbf{\ru{общество}} was taken from Church Slavonic \emph{obĭščĭstvo} (EtymDB), built on \emph{\ru{общий}} \textquotedblleft{}common, shared\textquotedblright{}: society is what people have in common.
\end{cousinweb}

\begin{grammarlens}[title={What you've built}]
One more word for telling, explaining and describing at B1.
\end{grammarlens}

\subsection*{Your turn}
\begin{itemize}
\raggedright
  \item \emph{Say it:} \emph{óbshchestvo}
  \item \emph{Say it:} \emph{óbshchestvo}, once more
  \item \emph{Say it:} a sentence of your own with \emph{óbshchestvo}
  \item \emph{Recall:} say the Russian for a competitor, then the Russian for to cooperate, then say \emph{óbshchestvo} again
\end{itemize}

\begin{tcolorbox}[breakable,colback=teal!4,colframe=teal!35!black,title={Before you move on}]
What does \ru{общество} mean? (\textquotedblleft{}Society\textquotedblright{}.) Say it once more.
\end{tcolorbox}
\section[\texorpdfstring{\ru{население} (naseléniye)}{население (naseléniye)}]{\ru{население} (naseléniye) — the population}

\label{lesson:RU-C364-naselenie}



\begin{quote}
Before the new one: say the Russian for to cooperate, then the Russian for society.
\end{quote}

\subsection*{You'll want to know: \ru{население}}
\textbf{\ru{население}} (\emph{naseléniye}) — \textquotedblleft{}the population\textquotedblright{}.

\begin{quote}
\ru{Население} \ru{города} — \ru{больше} \ru{миллиона} \ru{человек}.
\end{quote}

\begin{quote}
\emph{Naseléniye góroda — bólshe millióna chelovék.}
\end{quote}

\begin{quote}
The population of the town is over a million people.
\end{quote}

Neuter. \emph{\ru{Население} \ru{страны}} — a country’s population.

\begin{cousinweb}
\textbf{\ru{население}} is built on \emph{\ru{населить}} \textquotedblleft{}to people, to settle\textquotedblright{}, from \emph{\ru{на}-} and the root of \emph{\ru{селить}} \textquotedblleft{}to settle\textquotedblright{} and \emph{\ru{село}} \textquotedblleft{}a village\textquotedblright{}: the people settled in a place.
\end{cousinweb}

\begin{grammarlens}[title={What you've built}]
One more word for telling, explaining and describing at B1.
\end{grammarlens}

\subsection*{Your turn}
\begin{itemize}
\raggedright
  \item \emph{Say it:} \emph{naseléniye}
  \item \emph{Say it:} \emph{naseléniye}, once more
  \item \emph{Say it:} a sentence of your own with \emph{naseléniye}
  \item \emph{Recall:} say the Russian for to cooperate, then the Russian for society, then say \emph{naseléniye} again
\end{itemize}

\begin{tcolorbox}[breakable,colback=teal!4,colframe=teal!35!black,title={Before you move on}]
What does \ru{население} mean? (\textquotedblleft{}The population\textquotedblright{}.) Say it once more.
\end{tcolorbox}
\section[\texorpdfstring{\ru{общественный} (obshchéstvennyy)}{общественный (obshchéstvennyy)}]{\ru{общественный} (obshchéstvennyy) — public, social}

\label{lesson:RU-C364-obshchestvennyy}



\begin{quote}
Before the new one: say the Russian for society, then the Russian for the population.
\end{quote}

\subsection*{You'll want to know: \ru{общественный}}
\textbf{\ru{общественный}} (\emph{obshchéstvennyy}) — \textquotedblleft{}public, social\textquotedblright{}.

\begin{quote}
\ru{Я} \ru{езжу} \ru{на} \ru{работу} \ru{на} \ru{общественном} \ru{транспорте}.
\end{quote}

\begin{quote}
\emph{Ya yézzhu na rabótu na obshchéstvennom tránsporte.}
\end{quote}

\begin{quote}
I go to work by public transport.
\end{quote}

Public opinion is \emph{\ru{общественное} \ru{мнение}}, and public transport is \emph{\ru{общественный} \ru{транспорт}}.

\begin{cousinweb}
\textbf{\ru{общественный}} is the adjective made from \emph{\ru{общество}} \textquotedblleft{}society\textquotedblright{}, which Russian took from Church Slavonic \emph{obĭščĭstvo} (EtymDB).
\end{cousinweb}

\begin{grammarlens}[title={What you've built}]
One more word for telling, explaining and describing at B1.
\end{grammarlens}

\subsection*{Your turn}
\begin{itemize}
\raggedright
  \item \emph{Say it:} \emph{obshchéstvennyy}
  \item \emph{Say it:} \emph{obshchéstvennyy}, once more
  \item \emph{Say it:} a sentence of your own with \emph{obshchéstvennyy}
  \item \emph{Recall:} say the Russian for society, then the Russian for the population, then say \emph{obshchéstvennyy} again
\end{itemize}

\begin{tcolorbox}[breakable,colback=teal!4,colframe=teal!35!black,title={Before you move on}]
What does \ru{общественный} mean? (\textquotedblleft{}Public, social\textquotedblright{}.) Say it once more.
\end{tcolorbox}
\section[\texorpdfstring{\ru{социальный} (sotsiálnyy)}{социальный (sotsiálnyy)}]{\ru{социальный} (sotsiálnyy) — social}

\label{lesson:RU-C364-socialnyy}



\begin{quote}
Before the new one: say the Russian for the population, then the Russian for public.
\end{quote}

\subsection*{You'll want to know: \ru{социальный}}
\textbf{\ru{социальный}} (\emph{sotsiálnyy}) — \textquotedblleft{}social\textquotedblright{}.

\begin{quote}
\ru{Он} \ru{занимается} \ru{социальными} \ru{проблемами}.
\end{quote}

\begin{quote}
\emph{On zanimáyetsya sotsiálnymi problémami.}
\end{quote}

\begin{quote}
He deals with social problems.
\end{quote}

\emph{\ru{Социальные} \ru{сети}} — social networks. \emph{\ru{Социальная} \ru{помощь}} — welfare.

\begin{cousinweb}
\textbf{\ru{социальный}} is usually traced, through western languages, to Latin \emph{sociālis} \textquotedblleft{}of companionship, of allies\textquotedblright{}, built on \emph{socius} \textquotedblleft{}a companion\textquotedblright{} (Lewis \& Short).
\end{cousinweb}

\begin{grammarlens}[title={What you've built}]
One more word for telling, explaining and describing at B1.
\end{grammarlens}

\subsection*{Your turn}
\begin{itemize}
\raggedright
  \item \emph{Say it:} \emph{sotsiálnyy}
  \item \emph{Say it:} \emph{sotsiálnyy}, once more
  \item \emph{Say it:} a sentence of your own with \emph{sotsiálnyy}
  \item \emph{Recall:} say the Russian for the population, then the Russian for public, then say \emph{sotsiálnyy} again
\end{itemize}

\begin{tcolorbox}[breakable,colback=teal!4,colframe=teal!35!black,title={Before you move on}]
What does \ru{социальный} mean? (\textquotedblleft{}Social\textquotedblright{}.) Say it once more.
\end{tcolorbox}
\section[\texorpdfstring{\ru{житель} (zhítel)}{житель (zhítel)}]{\ru{житель} (zhítel) — an inhabitant, a resident}

\label{lesson:RU-C364-zhitel}



\begin{quote}
Before the new one: say the Russian for public, then the Russian for social.
\end{quote}

\subsection*{You'll want to know: \ru{житель}}
\textbf{\ru{житель}} (\emph{zhítel}) — \textquotedblleft{}an inhabitant, a resident\textquotedblright{}.

\begin{quote}
\ru{Жители} \ru{деревни} \ru{против} \ru{новой} \ru{дороги}.
\end{quote}

\begin{quote}
\emph{Zhíteli derévni prótiv nóvoy dorógi.}
\end{quote}

\begin{quote}
The villagers are against the new road.
\end{quote}

Masculine; a woman is \emph{\ru{жительница}}. \emph{\ru{Жители} \ru{города}} — the people of a town.

\begin{cousinweb}
\textbf{\ru{житель}} is built on \emph{\ru{жить}} \textquotedblleft{}to live\textquotedblright{}, inherited from Proto-Slavic \emph{*žiti}, with the agent suffix \emph{-\ru{тель}}: one who lives there.
\end{cousinweb}

\begin{grammarlens}[title={What you've built}]
One more word for telling, explaining and describing at B1.
\end{grammarlens}

\subsection*{Your turn}
\begin{itemize}
\raggedright
  \item \emph{Say it:} \emph{zhítel}
  \item \emph{Say it:} \emph{zhítel}, once more
  \item \emph{Say it:} a sentence of your own with \emph{zhítel}
  \item \emph{Recall:} say the Russian for public, then the Russian for social, then say \emph{zhítel} again
\end{itemize}

\begin{tcolorbox}[breakable,colback=teal!4,colframe=teal!35!black,title={Before you move on}]
What does \ru{житель} mean? (\textquotedblleft{}An inhabitant, a resident\textquotedblright{}.) Say it once more.
\end{tcolorbox}
